% GENERATED FILE. Edit canonical lessons, then run npm run generate:books.
% canonical-source-hash: fnv1a64:77a901fea121bb9f
% canonical-lessons: SA-C191-saravah, SA-C191-kumbhah, SA-C191-ulukhalam, SA-C191-musalam, SA-C191-calani

\chapter{Earthen dish, Pitcher, Mortar, Pestle, Sieve}
\label{ch:sa-earthen-dish-pitcher-mortar-pestle-sieve}
\hlchaptermodality{191}

\begin{chapteropening}
\textbf{By the end of this chapter:} \emph{I can say an earthen dish, a pitcher, a mortar, a pestle and a sieve.}

Complete the last lesson of chapter 191: I can say an earthen dish, a pitcher, a mortar, a pestle and a sieve.
\end{chapteropening}

\section[\texorpdfstring{\sk{शरावः} (śarāvaḥ)}{शरावः (śarāvaḥ)}]{\sk{शरावः} (śarāvaḥ) — an earthen dish}

\label{lesson:SA-C191-saravah}



\begin{quote}
Before the new one: say the Sanskrit for an amla, gooseberry, then the Sanskrit for a rose apple.
\end{quote}

\subsection*{You'll want to know: \sk{शरावः}}
\textbf{\sk{शरावः}} — \emph{śarāvaḥ} — \textquotedblleft{}an earthen dish\textquotedblright{}.

\begin{grammarlens}[title={What you've built}]
One more word for this chapter.
\end{grammarlens}

\subsection*{Your turn}
\begin{itemize}
\raggedright
  \item \emph{Say it:} \emph{śarāvaḥ}
  \item \emph{Say it:} \emph{śarāvaḥ}, once more
  \item \emph{Say it:} say \emph{śarāvaḥ}
  \item \emph{Recall:} say the Sanskrit for an amla, gooseberry, then the Sanskrit for a rose apple, then say \emph{śarāvaḥ} again
\end{itemize}

\begin{tcolorbox}[breakable,colback=teal!4,colframe=teal!35!black,title={Before you move on}]
What does \sk{शरावः} mean? (\textquotedblleft{}An earthen dish\textquotedblright{}.) Say it once more.
\end{tcolorbox}
\section[\texorpdfstring{\sk{कुम्भः} (kumbhaḥ)}{कुम्भः (kumbhaḥ)}]{\sk{कुम्भः} (kumbhaḥ) — a pitcher}

\label{lesson:SA-C191-kumbhah}



\begin{quote}
Before the new one: say the Sanskrit for a rose apple, then the Sanskrit for an earthen dish.
\end{quote}

\subsection*{You'll want to know: \sk{कुम्भः}}
\textbf{\sk{कुम्भः}} — \emph{kumbhaḥ} — \textquotedblleft{}a pitcher\textquotedblright{}.

\begin{grammarlens}[title={What you've built}]
One more word for this chapter.
\end{grammarlens}

\subsection*{Your turn}
\begin{itemize}
\raggedright
  \item \emph{Say it:} \emph{kumbhaḥ}
  \item \emph{Say it:} \emph{kumbhaḥ}, once more
  \item \emph{Say it:} say \emph{kumbhaḥ}
  \item \emph{Recall:} say the Sanskrit for a rose apple, then the Sanskrit for an earthen dish, then say \emph{kumbhaḥ} again
\end{itemize}

\begin{tcolorbox}[breakable,colback=teal!4,colframe=teal!35!black,title={Before you move on}]
What does \sk{कुम्भः} mean? (\textquotedblleft{}A pitcher\textquotedblright{}.) Say it once more.
\end{tcolorbox}
\section[\texorpdfstring{\sk{उलूखलम्} (ulūkhalam)}{उलूखलम् (ulūkhalam)}]{\sk{उलूखलम्} (ulūkhalam) — a mortar}

\label{lesson:SA-C191-ulukhalam}



\begin{quote}
Before the new one: say the Sanskrit for an earthen dish, then the Sanskrit for a pitcher.
\end{quote}

\subsection*{You'll want to know: \sk{उलूखलम्}}
\textbf{\sk{उलूखलम्}} — \emph{ulūkhalam} — \textquotedblleft{}a mortar\textquotedblright{}.

\begin{grammarlens}[title={What you've built}]
One more word for this chapter.
\end{grammarlens}

\subsection*{Your turn}
\begin{itemize}
\raggedright
  \item \emph{Say it:} \emph{ulūkhalam}
  \item \emph{Say it:} \emph{ulūkhalam}, once more
  \item \emph{Say it:} say \emph{ulūkhalam}
  \item \emph{Recall:} say the Sanskrit for an earthen dish, then the Sanskrit for a pitcher, then say \emph{ulūkhalam} again
\end{itemize}

\begin{tcolorbox}[breakable,colback=teal!4,colframe=teal!35!black,title={Before you move on}]
What does \sk{उलूखलम्} mean? (\textquotedblleft{}A mortar\textquotedblright{}.) Say it once more.
\end{tcolorbox}
\section[\texorpdfstring{\sk{मुसलम्} (musalam)}{मुसलम् (musalam)}]{\sk{मुसलम्} (musalam) — a pestle}

\label{lesson:SA-C191-musalam}



\begin{quote}
Before the new one: say the Sanskrit for a pitcher, then the Sanskrit for a mortar.
\end{quote}

\subsection*{You'll want to know: \sk{मुसलम्}}
\textbf{\sk{मुसलम्}} — \emph{musalam} — \textquotedblleft{}a pestle\textquotedblright{}.

\begin{grammarlens}[title={What you've built}]
One more word for this chapter.
\end{grammarlens}

\subsection*{Your turn}
\begin{itemize}
\raggedright
  \item \emph{Say it:} \emph{musalam}
  \item \emph{Say it:} \emph{musalam}, once more
  \item \emph{Say it:} say \emph{musalam}
  \item \emph{Recall:} say the Sanskrit for a pitcher, then the Sanskrit for a mortar, then say \emph{musalam} again
\end{itemize}

\begin{tcolorbox}[breakable,colback=teal!4,colframe=teal!35!black,title={Before you move on}]
What does \sk{मुसलम्} mean? (\textquotedblleft{}A pestle\textquotedblright{}.) Say it once more.
\end{tcolorbox}
\section[\texorpdfstring{\sk{चालनी} (cālanī)}{चालनी (cālanī)}]{\sk{चालनी} (cālanī) — a sieve}

\label{lesson:SA-C191-calani}



\begin{quote}
Before the new one: say the Sanskrit for a mortar, then the Sanskrit for a pestle.
\end{quote}

\subsection*{You'll want to know: \sk{चालनी}}
\textbf{\sk{चालनी}} — \emph{cālanī} — \textquotedblleft{}a sieve\textquotedblright{}.

\begin{grammarlens}[title={What you've built}]
One more word for this chapter.
\end{grammarlens}

\subsection*{Your turn}
\begin{itemize}
\raggedright
  \item \emph{Say it:} \emph{cālanī}
  \item \emph{Say it:} \emph{cālanī}, once more
  \item \emph{Say it:} say \emph{cālanī}
  \item \emph{Recall:} say the Sanskrit for a mortar, then the Sanskrit for a pestle, then say \emph{cālanī} again
\end{itemize}

\begin{tcolorbox}[breakable,colback=teal!4,colframe=teal!35!black,title={Before you move on}]
What does \sk{चालनी} mean? (\textquotedblleft{}A sieve\textquotedblright{}.) Say it once more.
\end{tcolorbox}
